\documentclass[11pt]{article}
\usepackage{graphicx}
\usepackage{amssymb}
\usepackage{bm}
\usepackage{mathtools}
\usepackage{amsmath}
\usepackage{commath}
\usepackage{amsthm}
\usepackage{enumitem}
\usepackage{float} 
\usepackage{array}
\usepackage{outlines}
\usepackage{setspace}
\usepackage[colorinlistoftodos]{todonotes}
\setlist[enumerate,1]{label=\alph*)}
\setlist[enumerate,2]{label=\roman*)}
\newtheorem{theorem}{Theorem}
\newtheorem{lemma}{Lemma}
\newtheorem{remark}{Remark}
\newtheorem{definition}{Definition}
\newtheorem{prop}{Proposition}
\newcommand*\diff{\mathop{}\!\mathrm{d}}
\usepackage{tikz}
\usetikzlibrary{arrows,shapes,trees, calc}

%%% *** PREAMBLE *** %%%
\title{Ideas and structure for the contributed talk.}

\date{}

%%%***  CONTENT *** %%%
\begin{document}

\maketitle

\section{The case of the $\epsilon$-free diffusion coefficients.}

Describe the model, the scaling and main result obtained for the $\epsilon$-free $\mathbf{M}$ diffusion matrix --- essentially what is described in the second submitted version of the paper.
  
\section{Discussing the meaning and the challenges of introducing $\epsilon$ to $K_1$}

\subsection{The physical idea: downwind-matching coordinate system.}
  
  %A widely used approach to mathematically describe the atmosphere is to consider it as a geophysical fluid in a shallow domain --- and thus to model it using classical fluid dynamical equations combined with the explicit inclusion of an $\epsilon$ parameter representing the small aspect ratio of the physical domain. In our previous paper we proved a weak convergence theorem for the polluted atmosphere described by the Navier-Stokes equations extended by an advection-diffusion equation. We obtained a justification of the generalised hydrostatic limit model including the pollution effect described in the case of a classical, east-north-upwards oriented local Cartesian coordinates.
  
  Here we give an improvement of the previous statement: we consider a meteorologically more meaningful coordinate system, incorporate the analytical consequences of this coordinate change into the governing equations, and verify that the weak convergence still holds for this altered concentration -- velocity system.
  
  A local domain on the surface of the rotating Earth is often described by a Cartesian coordinate system where the $x, y, z$ axes represent the directions towards east, north, and upwards, respectively, as in this specific system the rotation vector, and as a consequence, the Coriolis force takes a simpler form. However, for modelling wind and pollution effects in the atmosphere there is another naturally distinctive viewpoint for fixing the orientation of the coordinate system: the introduction of downwind and crosswind, and to match the horizontal axes according to these meteorological parameters. We take this idea as a foundation, and we migrate the result of our previously stated weak convergence theorem to this more considerately chosen coordinate system.

As in this scenario the advection is dominant compared to diffusion effects along the $x$ axis, this approach naturally brings an additional $\epsilon$ term in the Laplacian of the concentration equation, and brings the challenge of proving a weak convergence result for the product of the vertical velocity and the concentration, as in the energy inequality the $\norm{\partial_1 C^\epsilon}$ term is replaced by $\norm{\sqrt{\epsilon} \partial_1 C^\epsilon}.$ This issue is overcome by considering simplified boundary conditions that are easier to manage, and taking a more regular, $H^1$ initial horizontal velocity $v_0.$

As we highlighted before, the key point of this new model is to integrate some of the main ideas of a Gaussian plume dispersion model into the system describing the polluted atmosphere. Therefore we begin by fixing the local Cartesian coordinate system to be adjusted to the downwind direction, i.e. $z$ is oriented upwards, while $x$ is the downwind, and $y$ is the crosswind direction. Another assumption the Gaussian dispersion models typically use as well is that advection is more dominant than diffusion along the downwind direction, which leads to the appearance of an additional $\epsilon$ term in the concentration equation's Laplacian. Namely, we have

\begin{equation} \abs{u_1 \partial_{1} C} \gg \abs{ K_x \partial_{11} ^2 C }\end{equation}

which leads us to using this additional scaling equation

\begin{equation} \label{ksepsk1} K_x = \epsilon K_1.\end{equation}

\begin{remark} \normalfont
  The inclusion of these two features in a pollution model are legitimate when the following circumstances hold:
  \begin{itemize}
    \item We are considering a time interval which is not too long in the sense that it is reasonable to assume that we have a permanent, relatively stable main wind direction, i.e. fixing the horizontal coordinates is reasonable in the first place,
    \item There is an adequate, relatively high wind speed --- low wind speeds in general lead to a situation where three dimensional diffusion dominates, leading to an error caused by the dominant advection assumption. Fast changing wind direction and low windspeed resulting in fast pollutant settling can make it extremely challenging to obtain reliable results from Gaussian models.
  \end{itemize}
\end{remark}

Choosing the coordinate system in this particular way results in changes in three main points of the system. One of these is the Coriolis term, since now the $x$ axis is not necessarily perpendicular to the rotation vector. We need to revisit the boundary conditions as well, and as discussed above, a modification in the pollution concentration equation also has to be taken into account. We will describe these changes and their consequences in detail --- the main challenge in proving this scenario's convergence result is in fact to handle the effect of this latter change on the energy inequality.

In order to obtain the set of equations that describe the phenomena we start with the original set of equations

\begin{equation} \label{originalEQ-V}
  \partial_t \mathbf{v} + (\mathbf{v} \cdot \nabla)\mathbf{v} -\Delta_{\nu} \mathbf{v} + \nabla q + 2 \mathbf{\omega} \times \mathbf{v}= \mathbf{g} \text{ in } \Omega_{\epsilon} \times (0,T)
\end{equation}
    
\begin{equation} \label{originalEQ-I}
  \nabla \cdot \mathbf{v} = 0 \text{ in } \Omega_{\epsilon} \times (0,T)
\end{equation}
    
\begin{equation} \label{diffusive}
  \partial_t P + \mathbf{v} \cdot \nabla P = \nabla \cdot (\mathbf{K} \nabla P) + Q,
\end{equation}

As for the rescaling process, here we perform the same steps, with the addition of $(\ref{ksepsk1}).$

To expand the $2 \mathbf{\omega} \times \mathbf{v}$ Coriolis term firstly it is necessary to understand the structure of the new rotation vector. Note that the downwind-matching coordinate system and the original east-north-upwards oriented system are different only in a rotation around the $z$ axis. Let $\theta$ denote the angle between the respective $x$ axes, $\phi$ denotes the latitude, then the rotation vector expressed in the downwind matching coordinate system becomes

\begin{gather}
 \begin{bmatrix} \cos(\theta) & -\sin(\theta) & 0 \\ \sin(\theta) & \cos(\theta) & 0 \\ 0 & 0 & 1 \end{bmatrix} \begin{bmatrix} 0 \\ f \cos(\phi) \\ f \sin(\phi) \end{bmatrix}
 =
  f \begin{bmatrix} -\sin(\theta)\cos(\phi) \\ \cos(\theta) \cos(\phi) \\ \sin(\phi) \end{bmatrix}.
\end{gather}

Now, let

$$\alpha = -2 f \sin(\theta)\cos(\phi), \quad \beta = 2 f \cos(\theta) \cos(\phi), \quad \gamma = 2 f \sin(\phi),$$

and using this, the Coriolis force takes the form

\begin{gather}
\begin{bmatrix} \alpha \\ \beta \\ \gamma \end{bmatrix}
 \times
 \begin{bmatrix} u_1 \\ u_2 \\ \epsilon u_3 \end{bmatrix}
 =
 \begin{bmatrix}  \beta \epsilon u_3 - \gamma u_2 \\ \alpha \epsilon u_3 - \gamma u_1 \\ \alpha u_2 - \beta u_1 \end{bmatrix}.
\end{gather}

Thus, these are the terms that appear in the three velocity equations; and finally, note that the third component is present in the rescaled version of the vertical velocity equation multiplied by $\epsilon,$ which is obtained from the $\epsilon$ term in the denominator of $\partial_z p.$

\begin{remark} \normalfont
The intuitive meaning of what we understand by downwind and the velocity in this new coordinate system. The downwind direction is an approximately steady-in-time direction, it can be seen as the direction marked by the wind direction trackers on an airport or meteorological point. When we consider the downwind, we neglect the small perturbations in its orientation, it can be viewed as a time-averaged direction as long as the fluctuation is relatively small. On the other hand, the velocity is a three dimensional, time dependent vector that changes second after second. 

\end{remark}
    
    We arrive to the merged anisotropic equations of the polluted atmosphere above inland surface:
  
    \begin{equation} \label{ANSC1}
    \partial_t u^\epsilon _1 + \mathbf{u}^\epsilon \cdot \nabla u^\epsilon _1 -\Delta_\nu u^\epsilon _1 -\gamma u_2^\epsilon + \epsilon \beta u_3^\epsilon + \partial_1 p^\epsilon = 0 \text{ in } \Omega \times (0,T)
    \end{equation}
    
    \begin{equation} \label{ANSC2}
    \partial_t u_2^\epsilon + \mathbf{u}^\epsilon \cdot \nabla u_2^\epsilon -\Delta_\nu u_2^\epsilon + \alpha \epsilon u^\epsilon_3 - \gamma u_1^\epsilon + \partial_2 p^\epsilon = 0 \text{ in } \Omega \times (0,T)
    \end{equation}
    
    \begin{equation} \label{ANSC3}
    \epsilon^2 \{  \partial_t u_3^\epsilon + \mathbf{u}^\epsilon \cdot \nabla u_3^\epsilon -\Delta_\nu u_3^{\epsilon} \} + \epsilon \alpha u^\epsilon_2 - \epsilon \beta u_1^\epsilon + \partial_3 p^\epsilon = 0 \text{ in } \Omega \times (0,T)
    \end{equation}
    
    \begin{equation} \label{ANSC4}
    \nabla \cdot \mathbf{u}^\epsilon= 0 \text{ in } \Omega \times (0,T)
    \end{equation}
    
    \begin{equation} \label{ANSC5}
    \partial_t C^\epsilon + \mathbf{u^\epsilon} \cdot \nabla C^\epsilon = \epsilon K_1 \partial_{11} ^ 2 C^\epsilon + K_2 \partial_{22} ^ 2 C^\epsilon + K_3 \partial_{33} ^ 2 C^\epsilon + S_\epsilon \text{ in } \Omega \times (0,T)
    \end{equation}

The initial condition for the velocity and the pollution concentration are
    
    \begin{equation} \label{ANSC7}
    \mathbf{u}^\epsilon (\cdot, t = 0) = \mathbf{u}_0, \quad C^\epsilon (\cdot, t=0) = C_0  \quad \text{ in } \Omega.
    \end{equation}
  
    
If we assume $\mathbf{u}^\epsilon = O(1)$, then neglecting the $\epsilon^2$ and $\epsilon$ in the anisotropic equations, we formally arrive to the hydrostatic Navier-Stokes equations (i.e., primitive equations) combined with pollution.

    \begin{equation} \label{HNSC1}    
    \partial_t u_1 + \mathbf{u} \cdot \nabla u_1 - \Delta_\nu u_1 -\gamma u_2 + \partial_1 p = 0 \text{ in } \Omega \times (0,T)
    \end{equation}
    
    \begin{equation} \label{HNSC2}   
    \partial_t u_2 + \mathbf{u} \cdot \nabla u_2 - \Delta_\nu u_2 -\gamma u_1 + \partial_2 p = 0 \text{ in } \Omega \times (0,T)
    \end{equation}
    
    \begin{equation} \label{HNSC3}   
    \partial_3 p = 0 \text{ in } \Omega \times (0,T)
    \end{equation}
    
    \begin{equation} \label{HNSC4}   
    \nabla \cdot \mathbf{u} = 0 \text{ in } \Omega \times (0,T)
    \end{equation}
    
    \begin{equation} \label{HNSC5}   
    \partial_t C + \mathbf{u} \cdot \nabla C = K_2 \partial_{22} ^ 2 C + K_3 \partial_{33} ^2 C + S \text{ in } \Omega \times (0,T)
    \end{equation}
    
    \begin{equation} \label{HNSC7}   
    u_i (\cdot, t=0) = u_{0 i}, C (\cdot, t=0) = C_0 \text{ in } \Omega, \text{$i$ = 1,2}
    \end{equation}
 
  
  \subsection{The challenge caused by the additional $\epsilon$ term resulting in weaker regularities for $C$}
  
  We need to pass to the limit with the nonlinear term. $C$ is no longer in $H^1,$ so the original perturbation theorem can not be applied.
  
  \subsection{Periodic boundary conditions and strong convergence of $u_3$}
  
  Using periodic boundary conditions and the result that $v_0 \in H^1 \implies$ strong convergence of $u_3$ in $L^2 L^2$ using the Strong Convergence Justification article's result.
  
  \subsection{Main result}
  
  CHANGE: Actually, since we have the strong convergence of $u_3,$ we do not need this lemma, we can just directly use the strong convergence of $u_3$ to prove the weak convergence of a product whose one term if weakly convergent, and the other is strongly convergent.
  
  Using the strong convergence of $u_3$ for the case of an $H^1$ regular, periodic $v_0$, we can prove that the uniformity part of the equicontinuity estimate holds.
  
  Let $\epsilon = \epsilon_n$ be a sequence that goes to zero, and we choose $h^n = u^{\epsilon}_3, g^n = C^{\epsilon},$ and $q_1 = q_2 = p_1 = p_2 = 2,$ with $ m=-1.$

\textit{CONDITION1.} To estimate the time derivative of $C^{\epsilon}$ we use (this part for the concentration function is the same as before)

$$  \frac{\partial C^{\epsilon}}{\partial t} = - \nabla \cdot (\mathbf{u}^{\epsilon} C^{\epsilon}) + \sqrt{\epsilon} K_1 \partial_1 \partial_1 \sqrt{\epsilon} C^{\epsilon} + K_2 \partial_2 \partial_2 C^{\epsilon} + K_3 \partial_3 \partial_3 C^{\epsilon} -S^{\epsilon}. $$

The Gaussian source term is smooth; we have $ \nabla \cdot (\mathbf{u}^{\epsilon} C^{\epsilon}) \in W_x^{-1,1}$ and the remaining terms are in $W_x^{-1,2}.$ Eventually, putting these results together, $\partial_t C^{\epsilon} \in L_t^1 W_x^{-1,1}.$
         
\textit{CONDITION2.} The equicontinuity condition

$$\norm{u_3^{\epsilon}(\cdot, t) - u_3^{\epsilon} (\cdot + \xi, t)}_{L^2(0,T; L^2(\Omega))} \to 0 \text{ as } \abs{\xi} \to 0 \text{ \textbf{uniformly in }} \mathbf{\epsilon} $$

can be verified as follows. The time arguments are the same, we show that the spatial shifts vanish uniformly using that we have $H^1$ space regularity for $u^{\epsilon}_3.$
We have the compact embedding $H^1 \hookrightarrow L^2,$ the weak convergence of $u^{\epsilon}_3$ to $u_3$ in $L^2(\Omega),$ and the $H^1_x$ regularity for $u_3^{\epsilon},$ so we have that in space $u^{\epsilon}_3 \to u_3$ strongly in $L^2 (\Omega).$

Now we show that for any $E > 0$ there is a $\delta > 0$ for which any $\abs{\xi} < \delta,$ we have $\norm{u_3^{\epsilon}(\cdot) - u_3^{\epsilon}(\cdot + \xi)}_{L^2_\Omega} < E,$ independently of $\epsilon = \epsilon_n.$

Let's now fix this $E$ value, and we will define $\delta (E)$ as $\delta = \min{\{\delta_1, ... , \delta_{N-1}, \delta_N\}},$ where the $\delta_n$ values are defined below.

We know from the strong convergence of $u_3^{\epsilon}$ that for this $E$ there is an $N > 0$ for which $\forall n \geq N$ we have $\norm{u_3^{\epsilon_n}-u_3}_{L^2} < E / 3.$ Let $\delta_N $ be the value for which we have  $\norm{u_3(\cdot) - u_3(\cdot + \xi) }_{L^2(\Omega)} < E / 3$ for $\abs{\xi} < \delta_N.$ With this we have that $ \norm{u_3^{\epsilon}(\cdot) - u_3^{\epsilon}(\cdot + \xi)}_{L^2_\Omega},$ written as

$$\norm{u_3^{\epsilon}(\cdot) - u_3(\cdot) + u_3(\cdot) - u_3(\cdot + \xi) + u_3(\cdot + \xi) - u_3^{\epsilon}(\cdot + \xi)}_{L^2_\Omega},$$ is estimated by $ E / 3 + E / 3 + E / 3. $

On the other hand, for the values $n = 1, ... , N-1$ we define $\delta_1, ..., \delta_{N-1}$ such that $\norm{u_3^{\epsilon_n}(\cdot) - u_3^{\epsilon_n}(\cdot + \xi)}_{L^2_\Omega} < E$ when $\abs{\xi} < \delta_n.$

Let's take $\delta = \min{\{\delta_1, ... , \delta_{N-1}, \delta_N\}}.$

\end{document}

%%% add http://elte.prompt.hu/sites/default/files/tananyagok/AtmosphericChemistry/ch10s03.html
%%% http://twister.caps.ou.edu/PM2000/Chapter7.pdf
